\scanpage{127}
\uline{Théorème} \unsure{Le déterminant} de Fredholm de $\unsure{u(\in L^1)|u\in L^1}$ \illegible{1}

\unsure{Ntn|Notn} $\mathcal L_0$\,; dire \inserted{\unsure{est-ce que}} \unsure{exactement} quelles
\unsure{parties|partie} du \deleted{\unsure{th}} théorème \unsure{--} sont connues.

\uline{Théorème}

1) Pour \deleted{\unsure{tout}} $0<p\leq\infty$, $\mathcal L^p(\mathcal H)$ est un idéal \inserted{bilatère}
dans $L(\mathcal H)$, stable pour l'involution $A\to A^*$. De
\deleted{\unsure{On a}} façon précise :
\[
\deleted{S_p}\|A\|_p=\deleted{S_p}\|A^*\|_p \tag*{(1)}
\]
\[
\begin{aligned}
\deleted{S_p}\|BA\|_p&\leq\|B\|\deleted{S_p}\|A\|_p
\qquad\text{pour } A\in\mathcal L^p(\mathcal H),B\in L(\mathcal H)\\
\|AB\|_p&\leq\|B\|\ \|A\|_p
\end{aligned}\tag*{(2)}
\]
Si $\infty\geq p\geq1$, $A\to\|A\|_p$ est une norme sur $\mathcal L^p(\mathcal H)$ \deleted{\unsure{(2)|(top.}}
i.e. on a
\[
\|A+B\|_p\leq\|A\|_p+\|B\|_p\quad\|\underset{\text{(1)}}{\lambda} A\|_p=|\lambda|\|A\|_p
\quad(1\leq p\leq\infty\tag*{(3)}
\]
\inserted{et $\mathcal L^p(\mathcal H)$ est complet pour cette norme}
Si $0<p\leq\infty$, alors l'expression $S_p(A)=\|A\|_p^p$ sur
$\mathcal L^p(A)$ satisfait à l'inégalité du triangle
\[
S_p(A+B)\leq S_p(A)+S_p(B)\qquad(0<p\leq1)\tag*{(4)}
\]
et muni de la distance, invariante par \unsure{translat.|transl.}
$S_p(A-B)$, $\mathcal L^p(A)$ est un espace vectoriel \unsure{top. log.|topolog.}
métrisable et complet, non localement convexe
si $p<1$, et $\dim\mathcal H=+\infty$.
\margin{\textcircled{2} \uline{Dual de $\mathcal L^p$}\,; \uline{réflexivité}}

\deleted{2}3) Si $A\in\mathcal L^p(\mathcal H)$ $(0<p<+\infty)$, alors la suite
des valeurs propres \inserted{$\alpha_i$} de $A$ est de puissance
$p^{\mathrm{ième}}$ sommable, et on a
\[
\sum|\alpha_i|^p\leq S_p(A)\tag*{(5)}
\]
\margin{généralisation aux normes de Schatten? \unsure{Ici plutôt|Ceci plutôt} en dernier \uline{Corollaire}}
$S_p(A)$ et $\|A\|_p$ sont des fonctions croissantes
sur l'ensemble des opérateurs hermitiens
positifs, i.e $H$ et $K$ hermitiens positifs, $H\leq K$,
implique $S_p(H)\leq S_p(K)$.

\deleted{\unsure{3|2}}\unsure{4|3}) Soient $0<p,q,r\leq+\infty$, avec $\frac{1}{p}+\frac{1}{q}\geq\frac{1}{r}$. Si \deleted{\unsure{donc}}
$A\in\mathcal L^p(\mathcal H)$, $B\in\mathcal L^q(\mathcal H)$, alors $AB\in\mathcal L^r(\mathcal H)$, et on a
\[
\|AB\|_r\leq\|A\|_p\|B\|_q\tag*{(6)}
\]
\margin{(1) .....}
